\chapter{总结与展望}

\section{工作总结}

本实验完成了一个基于 LSSS 矩阵的 CP-ABE 实验实现。系统以访问控制树作为策略输入，自动构造 LSSS 矩阵，并使用线性重构判断用户属性集合是否满足策略。围绕这一主线，项目实现了 Setup、KeyGen、Encrypt 和 Decrypt 四个核心流程，并通过命令行脚本和烟雾测试验证了成功解密和拒绝解密两类行为。同时，系统在同一 \verb|Project| 工程中接入了基于 \verb|py-ecc| 的 BLS12-381 真实 pairing 后端，使同一 LSSS-CP-ABE 主线能够在真实 $G_1/G_2/G_T$ 群元素和配对运算上运行。

从理论角度看，本实验把 CP-ABE 的访问控制问题转化为线性代数问题。访问树中的 AND、OR 和阈值门被统一表示为 LSSS 矩阵行，解密时不再依赖针对具体树结构的递归判定，而是求解 $M_I^T\omega=e_1$。这一设计有助于理解 Waters 类 CP-ABE 构造中“线性秘密共享”和“双线性对抵消”之间的关系。

这一结论与主要参考资料的脉络一致：Bethencourt 等人~[2]给出了 CP-ABE 的密文策略访问控制模型，Waters~[3]将表达能力推进到 LSSS 矩阵表示，Liu 等人~[4]进一步讨论如何从更直观的阈值访问树生成 LSSS 矩阵。本文的贡献不是提出新的密码方案，而是把这些资料中的核心机制整理成可运行、可检查、可复现的课程实验系统。

从工程角度看，项目使用 PowerShell 和 .NET 标准库完成基础实现，具有较低复现门槛。密钥、密文、策略和矩阵均以 JSON 形式保存，便于实验检查、调试和结果分析。真实 pairing 路线进一步通过 Python 后端调用真实配对库，在保持 CLI 使用方式基本一致的同时，把底层群元素从整数指数推进到 BLS12-381 曲线点和 $G_T$ 元素。对称加密部分采用 AES-CBC-HMAC 组合，体现了“先加密再认证”的基本工程思想。

在结果呈现上，本文将公式、代码和实验现象对应起来。例如，第二章解释 LSSS 重构条件，第三章说明加解密公式，第四章给出关键代码片段，第五章再用九个场景验证这些设计。这样的组织方式使算法实现、程序输出和访问策略语义保持一致。

\section{主要完成内容}

本项目已经完成的主要内容包括：

\begin{enumerate}
    \item 定义访问控制树 JSON 格式，支持属性叶子和阈值门。
    \item 实现访问控制树到 LSSS 矩阵的自动转换。
    \item 实现模 $p$ 下的大整数运算、向量运算、快速幂和模逆。
    \item 实现 LSSS 重构系数求解，使用高斯消元判断属性集合是否满足策略。
    \item 实现 CP-ABE 的 Setup、KeyGen、Encrypt 和 Decrypt 流程。
    \item 实现 BLS12-381 真实 pairing 后端，支持 $G_1/G_2/G_T$ 群元素、哈希到 $G_1$、配对运算、序列化和基本成员检查。
    \item 使用 AES-CBC-HMAC 加密和认证实际明文载荷。
    \item 编写 CLI 脚本，使各阶段可以独立运行和观察结果。
    \item 编写烟雾测试，验证基础正反例、特殊属性名、密文篡改、$2$-of-$3$ 阈值门、重复属性和非法策略输入，并对真实 pairing 路线进行补充验证。
    \item 编写 LaTeX 报告，解释理论基础、方案设计、系统实现和实验结果。
\end{enumerate}

这些完成内容覆盖了课程题目中的两个关键词：一是“属性基加密方案”，对应 Setup、KeyGen、Encrypt 和 Decrypt 四个流程；二是“使用访问控制树构建 LSSS 矩阵”，对应策略 JSON、矩阵转换函数和重构系数求解函数。烟雾测试中的阈值门、重复属性和非法策略输入进一步说明系统不是只针对一个固定样例硬编码。

\section{实验价值}

本实验的价值主要体现在三个方面。

第一，它将抽象的访问结构转化过程具体化。LSSS 在论文中常以矩阵形式出现，但矩阵如何由策略树生成并不直观。本项目通过示例策略 $(A\wedge B)\vee C$ 展示了从树到矩阵的完整过程，使矩阵行、属性映射和重构系数之间的关系更加清楚。

第二，它把 CP-ABE 的正确性验证拆解为可观察的中间步骤。密文中保留了策略、矩阵、行属性和指数分量，使每个字段在解密中的作用都能与算法公式对应起来。

第三，它提供了由浅入深的实验平台。指数抽象模型保留代数结构，使课程重点集中在访问策略、秘密共享和算法流程上；真实 pairing 路线则把同一主线放入 BLS12-381 群元素和配对运算中运行，使实验从“公式结构可解释”进一步推进到“真实群运算可执行”。

除此之外，项目还具有一定的可扩展价值。由于策略解析、矩阵构造、重构求解和 CP-ABE 主流程之间的边界较清晰，真实 pairing 接入时主要替换对象是群运算和序列化层，访问树到 LSSS 的构造逻辑仍可保留。真实 pairing 接入结果说明，先建立指数抽象模型再替换底层群实现是一条可行的工程路线。

Riepel 和 Wee 的 FABEO 方案~[5]表明，真实 ABE 研究还会继续优化配对次数、密钥和密文大小，并讨论自适应安全和多密文安全等更强性质。与这些研究方向相比，本项目已经完成课程题目所要求的 LSSS-ABE 主线，但距离高性能、可部署的 ABE 系统仍有明显距离。

\section{不足与限制}

本项目仍存在明显限制，主要包括：

\begin{itemize}
    \item \textbf{安全证明仍未形式化完成。} 真实 pairing 路线已经使用 BLS12-381 真实群元素和配对运算，但还没有逐条证明其与 Waters 2011 安全模型完全一致。
    \item \textbf{指数抽象版仍只适合解释代数结构。} 该路线没有真实双线性群中的困难问题保证，主要用于快速复现和公式对应。
    \item \textbf{参数与库选择仍需进一步论证。} BLS12-381 后端提升了实现真实性，但正式安全声明仍需要说明曲线选择、哈希到群域分离、随机数采样和库版本假设。
    \item \textbf{策略输入形式较基础。} 当前需要手写 JSON 访问树，尚未支持自然布尔表达式解析。
    \item \textbf{输入校验仍有扩展空间。} 当前已能拒绝非法阈值策略和空属性，但对密文数组长度不一致等异常结构的检查还可以进一步细化。
    \item \textbf{测试覆盖仍有扩展空间。} 当前指数版与 pairing 版烟雾测试已覆盖核心正反例、一般阈值门、特殊属性名、密文篡改、重复属性和非法策略输入；复杂嵌套策略、畸形 JSON 和文件覆盖行为仍属于可继续补充的测试维度。
    \item \textbf{CLI 输出仍有扩展空间。} 当前已对属性名生成文件名时做安全化处理，但输出路径、错误信息和运行状态提示仍可以进一步规范。
\end{itemize}

这些限制构成本文安全性声明的边界。指数抽象模型能够说明公式结构和访问控制逻辑正确，但不覆盖真实部署环境中的困难问题假设；真实 pairing 路线已经覆盖真实群元素和配对调用，但仍需要更严格的安全模型对齐和实现级审查。对课程实验而言，这种边界说明表明项目区分了算法流程验证、真实群运算实现和形式化安全证明三个层次。

\section{真实安全证明条件对照}

为了进一步说明指数模型、真实 pairing 实现与论文安全证明之间的关系，表 \ref{tab:security-proof-conditions} 给出安全证明条件对照。这里以 Waters 2011 风格的 LSSS-CP-ABE 构造~[3]为主要参照，关注实现是否满足论文安全模型的前提。

\begingroup
\renewcommand{\arraystretch}{1.18}
\begin{longtable}{>{\raggedright\arraybackslash}p{0.18\textwidth} >{\raggedright\arraybackslash}p{0.74\textwidth}}
\caption{安全证明条件对照}
\label{tab:security-proof-conditions}\\
\toprule
条件 & 对照说明 \\
\midrule
\endfirsthead
\toprule
条件 & 对照说明 \\
\midrule
\endhead
\textbf{双线性群} &
\textbf{真实实现：}使用素数阶配对友好曲线，提供 $G_1,G_2,G_T$ 和非退化双线性映射 $e:G_1\times G_2\rightarrow G_T$。\newline
\textbf{当前状态：}指数抽象路线使用 \texttt{ExponentSimulation}；真实 pairing 路线使用 \texttt{py-ecc} 的 BLS12-381 $G_1/G_2/G_T$ 与 \texttt{pairing}，已具备真实配对运算基础。\\[0.45em]
\textbf{群运算接口} &
\textbf{真实实现：}群元素乘法、幂运算、逆元、配对运算和目标群运算均由真实密码库实现。\newline
\textbf{当前状态：}\texttt{-Adapter pairing} 路线已由 Python 后端实现真实群元素计算；指数路线继续保留为解释和对照实现。\\[0.45em]
\textbf{随机数生成} &
\textbf{真实实现：}所有 $\mathbb{Z}_p$ 标量必须由密码学安全随机源均匀采样。\newline
\textbf{当前状态：}指数版使用 .NET 随机数生成器；pairing 版使用 Python \texttt{secrets.randbelow} 在曲线阶内采样。正式安全声明仍需进一步论证采样和参数选择。\\[0.45em]
\textbf{属性映射} &
\textbf{真实实现：}属性应按论文要求映射到群元素或公共参数；若使用哈希到群，应采用标准哈希到曲线方法。\newline
\textbf{当前状态：}指数版把属性哈希为标量；pairing 版调用 \texttt{hash\_to\_G1}，并使用固定 DST 将属性字符串映射为 $G_1$ 点。\\[0.45em]
\textbf{序列化检查} &
\textbf{真实实现：}反序列化群元素时必须检查编码合法性、曲线成员关系、子群成员关系和无穷远点。\newline
\textbf{当前状态：}pairing 版 JSON 保存 $G_1/G_2$ 坐标和 $G_T$ 系数，并在反序列化时检查曲线成员关系和子群成员关系；仍可继续增强无穷远点拒绝策略和异常密文测试。\\[0.45em]
\textbf{算法一致性} &
\textbf{真实实现：}Setup、KeyGen、Encrypt、Decrypt 的群元素位置、指数关系和访问结构必须与选定论文方案一致。\newline
\textbf{当前状态：}四算法结构与本文采用的 LSSS-CP-ABE 实验公式一致；pairing 版已经将 $K,L,K_x,C,C_0,C_i,D_i$ 分别落到 $G_1/G_2/G_T$ 对应位置。\\[0.45em]
\textbf{安全声明范围} &
\textbf{真实实现：}只能声明与论文模型相匹配的安全性，例如选择安全或自适应安全，并说明 CPA/CCA 范围。\newline
\textbf{当前状态：}实验结论限定为 LSSS-CP-ABE 的代数正确性、访问控制语义和真实 pairing 路径可运行性；不扩展为完整真实部署安全性。\\
\bottomrule
\end{longtable}
\endgroup

基于上述对照，可以得到一个更准确的结论：当前项目已经从单纯指数抽象验证推进到真实 BLS12-381 pairing 初步实现，验证了 LSSS-CP-ABE 的代数正确性、访问控制语义和真实群运算路径可运行性；若要进一步满足论文级安全证明条件，还需要将 Setup、KeyGen、Encrypt、Decrypt、哈希到群、随机数生成、序列化检查和安全声明逐项与选定论文的安全模型严格对齐。

\section{后续工作展望}

当前实现已经覆盖课程题目要求的 LSSS-ABE 主线，并在同一 \verb|Project| 工程中完成真实 pairing 初步接入。后续工作的核心不是重新设计访问控制树或 LSSS 矩阵构造，而是在保持现有模块边界的基础上，把真实 pairing 实现的安全模型对齐、策略表达能力、测试覆盖范围和实验呈现一致性进一步推进。四个方向之间也存在依赖关系：真实配对后端决定系统能否进入更接近真实密码实现的层面，策略语言和输入校验决定外部策略表达的可靠性，测试集决定边界行为是否稳定，实验呈现则负责把这些变化反映到报告证据链中。

\subsection{完善真实双线性对库实现}

真实 pairing 路线已经使用 \verb|py-ecc| 接入 BLS12-381，提供 $G_1$、$G_2$、$G_T$、$\mathbb{Z}_p$ 和配对映射 $e:G_1\times G_2\rightarrow G_T$。这一步使系统不再只停留在整数指数模拟层面，而是能够生成真实曲线点、调用真实配对函数，并把 $G_T$ 元素用于会话密钥派生。后续完善的重点，是把这一可运行实现继续向论文安全模型和工程健壮性靠拢。

现有结果说明群运算适配器边界是有效的。访问树解析、LSSS 矩阵生成和重构系数求解都只依赖标量和矩阵，不依赖具体群元素编码；因此真实化过程中稳定性较高的部分是策略层和线性代数层，变化最大的部分是 Setup、KeyGen、Encrypt、Decrypt 中的群元素表示、序列化格式和配对计算。换言之，当前实现已经把“访问结构正确性”和“底层群实现真实性”分开，这使系统能够同时保留指数抽象版和真实 pairing 版。

真实 pairing 版仍有进一步加固空间。群元素反序列化已经检查曲线成员关系和子群成员关系，但还可以加入更细的无穷远点拒绝、字段范围检查和异常密文测试；属性哈希已经使用 \verb|hash_to_G1|，但仍需在报告中固定域分离标签和哈希套件说明；随机标量采样已经使用安全随机源，但正式安全声明仍需说明采样均匀性；对称密钥派生则需要明确 $G_T$ 元素编码到固定长度密钥的规范。这些问题不属于 LSSS 矩阵转换本身，却直接影响真实 CP-ABE 实现能否更严谨地满足论文安全模型。

以实验三课件~[1]和 Waters 2011 方案~[3]为参照，真实化版本的比较对象还包括 FABEO 方案~[5]。Waters 方案更适合验证 LSSS 矩阵和重构常数进入加解密公式的方式，FABEO 方案则体现了较新 ABE 研究对密文大小、私钥大小、配对次数和安全模型的进一步优化。这样的对比能够把本实验从单一实验原型推进到不同 CP-ABE 构造之间的结构比较。

\subsection{增强策略语言}

当前系统使用 JSON 访问树作为策略输入，内部结构清晰，适合程序直接读取和递归校验；不足之处在于策略书写形式较接近中间表示，不够接近自然权限表达。策略语言扩展的目标，是在 JSON 访问树之前增加一层布尔表达式解析，外部输入形式例如：

\[
    (A \text{ AND } B) \text{ OR } C.
\]

解析层的输出仍然是当前系统已经支持的访问树结构。由此形成外部语法和内部表示的解耦：外部语法面向策略书写，内部 JSON 面向矩阵构造。AND、OR 和阈值门仍然统一归约为 $k$-of-$n$ 节点，后续 LSSS 转换函数无需改变。对于 \verb|role:admin|、\verb|dept:research| 这类带命名空间的属性，解析器还承担区分属性名冒号与语法操作符的任务，避免把合法属性名误判为非法表达式。

策略语言增强也与输入校验相关。表达式解析阶段能够发现括号不匹配、连续操作符、空属性、未知门类型等语法错误；访问树校验阶段则继续负责检查阈值 $k$、子节点数量和叶子节点结构。两级校验分别对应“外部策略文本是否合法”和“内部访问树是否能生成 LSSS 矩阵”，使错误来源更加明确。

\subsection{扩展测试集}

现有烟雾测试已经覆盖基础正反例、特殊属性名、密文篡改、一般阈值门、重复属性和非法策略输入。进一步扩展测试集时，重点在于增加策略结构深度、失败路径种类和异常密文形态，使测试不只覆盖示例策略，也覆盖更复杂的 LSSS 构造情形。

\begin{table}[H]
\centering
\caption{后续测试覆盖维度}
\begin{tabular}{>{\raggedright\arraybackslash}p{0.22\textwidth} >{\raggedright\arraybackslash}p{0.32\textwidth} >{\raggedright\arraybackslash}p{0.34\textwidth}}
\toprule
测试维度 & 代表场景 & 分析价值 \\
\midrule
嵌套策略 & 多层 AND/OR/阈值门组合 & 验证递归份额传播和随机列分配是否稳定。\\
失败路径 & 只满足局部子树但不满足根节点 & 验证线性重构不会把局部满足误判为全局满足。\\
异常密文 & 缺失 \verb|Ci|、\verb|Di| 或数组长度不一致 & 验证解密阶段对密文字段结构的检查能力。\\
策略语法 & 畸形 JSON、非法阈值、空属性、重复属性 & 验证策略校验能在矩阵构造前截断错误输入。\\
文件行为 & 输出路径已存在、属性名含特殊字符 & 验证 CLI 层对工程边界的处理稳定性。\\
\bottomrule
\end{tabular}
\end{table}

这些测试维度分别覆盖算法正确性和工程健壮性。嵌套策略和失败路径更接近 LSSS 理论层面，能够暴露矩阵列数、行映射和重构系数求解中的问题；异常密文、策略语法和文件行为更接近工程层面，能够暴露 JSON 结构、路径处理和错误信息中的问题。二者结合后，测试结果才能较充分地支撑“系统不仅能跑通示例，也能处理边界情况”的结论。

\subsection{实验呈现的一致性}

本报告的实验呈现围绕“理论定义、程序实现、文件结果和运行现象”之间的对应关系展开。系统架构图用于说明模块边界，加解密数据流图用于说明四个核心算法之间的数据传递，策略到 LSSS 矩阵的转换图用于说明访问树如何进入线性秘密共享结构；关键代码节选、密文 JSON 字段、烟雾测试输出和真实终端截图则分别对应实现依据、文件证据和运行结果。由此，报告中的图、表、代码和截图不是相互独立的材料，而是共同指向同一条实验链路。

这种呈现方式使第五章的实验结果能够被追溯到前文的算法设计。成功解密场景不仅有终端输出，也能对应到访问策略、LSSS 行集合和重构方程；失败解密场景不仅表现为错误信息，也能对应到属性集合无法满足 $M_I^T\omega=e_1$ 的线性条件；密文篡改、重复属性和非法策略输入等边界场景，则分别对应对称载荷完整性、行属性映射 $\rho$ 和访问树校验逻辑。实验报告因此不是只展示“程序可以运行”，而是把运行现象解释为算法结构的结果。

从报告完整性看，当前材料已经覆盖课程实验所需的主要证据类型：公式给出算法依据，代码节选说明关键实现，JSON 文件展示中间结果，截图记录真实运行，测试表归纳实验现象。真实 pairing 版已经补充了底层群运算真实性证据；若进一步推进策略语言、异常输入测试和安全模型对齐，也应延续同样的呈现原则，即让新增实现能够在理论定义、代码结构、文件字段和实验输出之间形成可核对的对应关系。

\section{结论}

总体而言，本项目完成了一个以 LSSS 矩阵为核心的 CP-ABE 实验实现。系统从访问控制树出发，自动生成 LSSS 矩阵 $(M,\rho)$，并将矩阵行、用户属性集合和重构系数统一到同一个解密判定过程之中。与只演示一次加密和解密的简单程序相比，本项目更强调访问结构的可解释性：策略为什么满足、矩阵如何重构、属性不足时为什么失败，都可以在公式、代码和实验结果之间找到对应关系。

从课程实验要求看，本项目覆盖了“属性基加密方案”和“使用访问控制树构建 LSSS 矩阵”两个核心点。Setup、KeyGen、Encrypt 和 Decrypt 四个流程已经形成完整闭环；访问控制树、LSSS 矩阵、行属性映射、重构方程和密文载荷也都以可检查的 JSON 或报告图表形式呈现。第五章的九个测试场景进一步说明，系统不仅支持示例策略 $(A\wedge B)\vee C$，也能处理特殊属性名、密文篡改、一般阈值门、重复属性和非法策略输入等边界情况。

同时，本文对实现边界作出了明确区分。指数抽象版用于验证 LSSS-CP-ABE 的代数正确性和访问控制语义，帮助理解 Waters 类 CP-ABE 构造中线性秘密共享和配对抵消之间的关系；真实 pairing 路线使用 BLS12-381 群元素和配对函数，进一步说明底层群运算真实化后该主线仍可运行。若要进一步满足真实论文中的安全证明条件，还需要继续审查哈希到群、随机数生成、序列化检查、参数选择和安全模型声明。

因此，本项目的最终结论是：当前实现已经完成课程层面的 LSSS-ABE 实验目标，能够较清楚地展示从访问策略到矩阵、从矩阵到重构、从重构到解密的完整过程；真实 pairing 路线已经完成初步接入，使项目在实现层面从指数抽象验证推进到 BLS12-381 群运算验证。后续工作的重点不在于重新设计访问树和 LSSS 主线，而在于沿着现有模块边界继续完善安全模型对齐、策略语言表达能力和异常输入测试覆盖。这样既保留了本实验已经建立的可解释结构，也为更接近真实密码系统的实现留下了明确扩展路径。
